% !TeX root = practica_limites_asintotas.tex
% !TeX program = pdflatex
% Práctica de los contenidos de límites de las semanas 03 y 04.
\documentclass[10pt,a5paper]{article}
\newif\ifmostrarSoluciones
\mostrarSolucionestrue
\usepackage{estilo_apuntes}
\newtcolorbox[auto counter]{ejerciciopractica}[1][]{
  breakable,
  colback=GrisClaro,
  colframe=Azul,
  title={\titulonumerado{Ejercicio}{\thetcbcounter}{#1}},
  fonttitle=\bfseries,
  arc=2mm,
  boxrule=0.7pt
}
\newenvironment{solucionejercicio}[1]{%
  \par\medskip\noindent\textbf{Solución del ejercicio #1.}%
  \enspace\ignorespaces
}{%
  \par\medskip
}
\encabezado{Semana 5}{Práctica de límites}

\begin{document}
\tituloclase{Práctica de límites y asíntotas}{Límites infinitos y al infinito, racionalización y asíntotas}{Contenidos de las semanas 3 y 4}

\begin{objetivos}
\begin{itemize}
  \item Determinar límites infinitos mediante el análisis de signos e identificar asíntotas verticales.
  \item Calcular límites al infinito de funciones racionales y radicales, e identificar asíntotas horizontales.
  \item Justificar cada respuesta conservando las restricciones del dominio.
\end{itemize}
\end{objetivos}

\noindent\textbf{Indicaciones.} Muestre el desarrollo e indique las propiedades que utiliza. En los ejercicios con asíntotas, escriba la ecuación de cada recta y el límite que la justifica.

\begin{ejerciciopractica}[radicales y dos asíntotas horizontales]
  Sea
  \[
    H(x)=\frac{\sqrt{4x^2+9}}{x-1},
    \qquad D_H=\mathbb{R}\setminus\{1\}.
  \]
  \begin{enumerate}
    \item Calcule sus límites cuando $x\to+\infty$ y cuando $x\to-\infty$.
    \item Determine sus asíntotas horizontales y verticales, justificando cada una mediante límites.
  \end{enumerate}
\end{ejerciciopractica}

\ifmostrarSoluciones
\begin{solucionejercicio}{1}
  Para $x\neq0$,
  \[
    \sqrt{4x^2+9}=|x|\sqrt{4+9/x^2}.
  \]
  Si $x\to+\infty$, se tiene $|x|=x$, de modo que
  \[
    H(x)=\frac{\sqrt{4+9/x^2}}{1-1/x}\longrightarrow2.
  \]
  Si $x\to-\infty$, se tiene $|x|=-x$, por lo que
  \[
    H(x)=-\frac{\sqrt{4+9/x^2}}{1-1/x}\longrightarrow-2.
  \]
  Las asíntotas horizontales son $y=2$ hacia $+\infty$ y $y=-2$ hacia $-\infty$.

  Cerca de $1$, el numerador es positivo y tiende a $\sqrt{13}$. El denominador tiende a $0$ con el signo de $x-1$. Así,
  \[
    \lim_{x\to1^-}H(x)=-\infty,
    \qquad \lim_{x\to1^+}H(x)=+\infty.
  \]
  La única asíntota vertical es $x=1$: el radicando es siempre positivo y el denominador no se anula en ningún otro punto.
\end{solucionejercicio}
\fi

\begin{ejerciciopractica}[racionalización al infinito]
  Considere la función definida en $\mathbb{R}$ por
  \[
    J(x)=\sqrt{x^2+6x+13}-x.
  \]
  \begin{enumerate}
    \item Calcule $\lim_{x\to+\infty}J(x)$ mediante racionalización.
    \item Calcule $\lim_{x\to-\infty}J(x)$.
    \item Determine las asíntotas horizontales e indique en qué extremo se presentan.
  \end{enumerate}
\end{ejerciciopractica}

\ifmostrarSoluciones
\begin{solucionejercicio}{2}
  El radicando es $(x+3)^2+4>0$, así que la función está definida para todo número real.

  Cuando $x\to+\infty$, aparece una diferencia de términos que crecen sin cota. Para $x>0$, multiplicamos por el conjugado y obtenemos
  \begin{align*}
    J(x)&=\frac{6x+13}{\sqrt{x^2+6x+13}+x}\\
        &=\frac{6+13/x}{\sqrt{1+6/x+13/x^2}+1}.
  \end{align*}
  En el último paso usamos $\sqrt{x^2}=x$, válido porque $x>0$. Por tanto,
  \[
    \lim_{x\to+\infty}J(x)=\frac{6}{1+1}=3.
  \]

  Para $x<0$, la raíz es no negativa y $-x>0$, de modo que $J(x)\geq-x$. Como $-x\to+\infty$ cuando $x\to-\infty$, la comparación implica
  \[
    \lim_{x\to-\infty}J(x)=+\infty.
  \]
  La única asíntota horizontal es $y=3$, hacia $+\infty$. No hay asíntota horizontal hacia $-\infty$.
\end{solucionejercicio}
\fi

\begin{ejerciciopractica}[integración: parámetro, límites y asíntotas]
  Para $a\in\mathbb{R}$, considere
  \[
    F_a(x)=\frac{2x^2+ax-3}{x^2-1},
    \qquad D_{F_a}=\mathbb{R}\setminus\{-1,1\}.
  \]
  \begin{enumerate}
    \item Determine $a$ para que $\lim_{x\to1}F_a(x)$ sea finito.
    \item Para ese valor de $a$, calcule el límite anterior y determine todas las asíntotas verticales y horizontales de $F_a$.
  \end{enumerate}
\end{ejerciciopractica}

\ifmostrarSoluciones
\begin{solucionejercicio}{3}
  Cuando $x\to1$, el numerador tiende a $a-1$ y el denominador tiende a $0$. Si $a\neq1$, el numerador conserva un signo no nulo cerca de $1$, mientras que el denominador cambia de signo. En ese caso no hay límite finito. Es necesario que $a=1$; comprobemos que también es suficiente.

  Para $a=1$ se tiene
  \[
    2x^2+x-3=(x-1)(2x+3),\qquad
    x^2-1=(x-1)(x+1).
  \]
  Por tanto, para $x\neq-1,1$,
  \[
    F_1(x)=\frac{2x+3}{x+1}=2+\frac{1}{x+1}.
  \]
  En consecuencia,
  \[
    \lim_{x\to1}F_1(x)=\frac52.
  \]
  El punto $1$ sigue excluido del dominio original, pero no produce una asíntota vertical.

  En $-1$, la expresión $2+1/(x+1)$ permite obtener
  \[
    \lim_{x\to-1^-}F_1(x)=-\infty,
    \qquad\lim_{x\to-1^+}F_1(x)=+\infty.
  \]
  Así, la única asíntota vertical es $x=-1$. Finalmente,
  \[
    \lim_{x\to+\infty}F_1(x)=
    \lim_{x\to-\infty}F_1(x)=2,
  \]
  de modo que la asíntota horizontal es $y=2$ hacia ambos extremos.
\end{solucionejercicio}
\fi

\end{document}
